\documentclass[a4paper]{article}
\usepackage[pdftex]{graphicx}
\usepackage[utf8]{inputenc}
\usepackage{enumerate}
\usepackage{icomma}
\usepackage{amssymb}
\usepackage{tikz}
\usepackage{href-ul}
\hypersetup{
	colorlinks=true,
	linkcolor=blue,
	urlcolor=blue}
\usepackage{geometry}
\geometry{a4paper, top=15mm, left=15mm, right=15mm, bottom=15mm,
	headsep=10mm, footskip=12mm}

\begin{document}
\begin{tikzpicture}(10,3)
	\draw[ultra thick](2,0) --(10,0) -- (10,1.3) --(2,1.3) -- (2,0);
	\draw[fill=black](2,0)-- (0,.6) -- (2,.6) -- (2,0);
	\node at (6,.6) {\large Lösungsblatt von \href{https://www.okuyakl.de}{www.okuyakl.de}};
\end{tikzpicture}
\vspace{0.5 cm}

\noindent{\bf Aufgabe 1. a)}\\
Die gesuchte Funktion soll vierten Grades sein, also ist unsere Grundform:
$$f(x)=ax^4+bx^3+cx^2+dx+e$$
Eine ganzrationale Funktion ist achsensymmetrisch, wenn sie nur gerade Potenzen von x hat. Dadurch reduziert sich unsere Form auf:
$$f(x)=ax^4+cx^2+e$$
Nun setzen wir die Nullstelle ein und erhalten:
$$f(0)= 0  \quad \Rightarrow \quad e=0$$
Weiter setzen wir den Hochpunkt ein:
$$f(1)=2 \quad \Rightarrow \quad a + c = 1 \quad I$$ 
Im Hochpunkt ist die Ableitung null:
$$f'(x)=4ax^3 + 2cx \stackrel{!}{=}0  \quad \Rightarrow \quad f'(1)= 4a + 2c =0 \quad II$$
Wir lösen I nach a auf und setzen in II ein:
$$4(1-c) + 2c =0 \quad \Rightarrow \quad c= c=2 $$
Daraus folgt $a=-1$
Der Funktionsterm ist also:
$$f(x)=-x^4+2x^2$$

\noindent{\bf Aufgabe 1. b)}\\
Zunächst benötigen wir die Nullstellen rechts und links:
$$f(x)=-x^4+2x^2=0 \quad \Rightarrow \quad x_{1/2}=\pm\sqrt{2}$$
Die Fläche ist gleich dem Integral dazwischen:
$$\int\limits_{-\sqrt{2}}^{\sqrt{2}}(-x^4+2x^2) {\rm d} x =\left[-{1\over 5}x^5+{2 \over 3}x^3\right]_{-\sqrt{2}}^{\sqrt{2}} = {16\over 15}\sqrt{2} \approx 1,51~{\rm FE} $$ 

\noindent{\bf Aufgabe 1. c)}\\
Wir subtrahieren die Funktion von der Geraden $y=1$ und integrieren:
$$\int\limits_{-1}^1 (1-(-x^4+2x^2)){\rm d} x =\left[x+{1\over 5}x^5-{2 \over 3}x^3\right]_{-1}^1={16 \over 15}~{\rm FE} $$


\begin{minipage}{0.5\textwidth}
	\noindent {\bf Aufgabe 2. a)}\\
    \includegraphics[width=6 cm]{exuf063}
\end{minipage}
\hfill
\begin{minipage}{0.5\textwidth}
 Der Graph ist gegenüber der ursprünglichen e-Funktion um zwei nach unten und eins nach rechts verschoben.\\
 \noindent {\bf Aufgabe 2. b)}\\
 $$H(x)=e^{x-1}-2x$$
\end{minipage}

\noindent {\bf Aufgabe 3. a)}\\
Wir leiten ab mit der Produktregel: 
$$F'(x)= e^{-x}\cdot(x+1)-e^{-x} = x\cdot e^{-x}=f(x)$$

\noindent {\bf Aufgabe 3. b)}\\
Wir integrieren von 0 bis $z$:
$$ A(z)=\int\limits_{0}^z f(x){\rm d} x=\left[ -e^{-x}\cdot(x+1)\right]_0^z=-e^{-z}\cdot(z+1)+1$$
Anschließend bilden wir den Grenzwert:
$$ A = \lim\limits_{z\to \infty} -e^{-z}\cdot(z+1)+1 = 1$$
Die Fläche hat den endlichen Wert 1 FE.

\noindent {\bf Aufgabe 4. a)}\\
Das Integral ist gleich:
$$
\renewcommand{\arraystretch}{2}
\begin{array}{rcccl}
\left[ 4 \sin x \right]_0^a =& 4 \sin a &=& 2 &|:4\\
 &  \sin a &=& {1\over 2}  \\
 &  a_1 &=& {1 \over 6}\pi \\
 &  a_2 &=& {5 \over 6}\pi
\end{array}
$$   

\noindent {\bf Aufgabe 4. b)}\\
Das Integral ist gleich:
$$\left[ {1\over 2} (-\cos x )\right]_0^2 = 0,827$$
Die Funktion $f(x)=\sin(2x)$ hat eine Nullstelle bei ${\pi \over 2}\approx 1,5$. Wir haben im Intervall $[0;2]$ ein Flächenstück oberhalb der x-Achse und eins unterhalb, welches negativ in die Flächenbilanz eingeht.

\noindent {\bf Aufgabe 5. a)}\\
Die Parallele zur x-Achse hat die Funktionsgleichung $y=k^2$.
Der Wert des Integrals als Flächenbilanz soll dann null werden:
$$\int\limits_0^1 (x^2 - k^2) {\rm d} x =0$$
damit folgt für $k$:
$$
\renewcommand{\arraystretch}{2}
\begin{array}{rcll}
\left[{1\over3 }x^3 - k^2 x \right]_0^1 &=& 0 \\
{ 1\over 3} - k^2 &=& 0 \\
k &=& \pm \sqrt{1 \over 3} \\ 
\end{array}
$$ 
Es kommt nur die positive Lösung infrage, also ist $k={1\over 3} \sqrt{3}$.

\noindent{\bf Aufgabe 6. a)}\\

Nullstellen: 
$$
\begin{array}{rcl}
-x^2+3x&=&0 \\
x\cdot(-x+3)&=&0 \quad x_1=0 \\
x_2&=&3 \Rightarrow N_1(0|0); \quad N_2(3|0)
\end{array}
$$


\noindent
\begin{minipage}{0.6\textwidth}
	Um auf Extrempunkte zu untersuchen bilden wir die erste Ableitung :
	
	$$ f'(x)=-2x+3 \stackrel{!}{=}0 \quad -2(x-1,5)=0 \quad \Rightarrow x_0 = 1,5 $$
	Vorzeichentabelle zur Bestimmung des Extrempunktes (hier nicht verlangt, aber für später folgende Aufgabe relevant):
	$$
	\begin{array}{cc|c|c}
	& x<1,5 & 1,5 & x>1,5\\
	\hline
	-2 & - & - & - \\
	(x-1,5) & - & 0 & + \\
	\hline
	f'(x) & + & 0 & - \\
	\hline
	f(x) & \nearrow & \rightarrow & \searrow \\
	& & HP & 
	\end{array}
	$$
	
	Die y- Koordinate des Hochpunktes ist:
	$$ f(1,5) = -1,5^2+3 \cdot 1,5 = 2,25 \quad \Rightarrow H(1,5|2,25)$$
\end{minipage}
\hfill
\begin{minipage}{0.4\textwidth}
\includegraphics[width=6 cm]{Diagramm140}
\end{minipage}




\noindent {\bf Aufgabe 6. b)}

Flächeninhalt:

$$A_{ges}=\int \limits_0^3 (-x^2+3x) {\rm d} x = \left[ -{1 \over 3} x^3 + {3\over 2} x^2 \right]_0^3= {9 \over 2} $$

\noindent {\bf Aufgabe 6. c)}

Schnittpunkt von Gerade und Parabel:

$$ 
\begin{array}{rcl}
-x^2+3x & = & x \\
x^2 &=&  2x \Rightarrow x_1=0; \quad x_2=2 
\end{array}
 $$

Flächeninhalt zwischen Gerade und $G_f$:

$$A_1=\int \limits_0^2 (-x^2+3x-x) {\rm d} x = \left[ -{1 \over 3} x^3 +  x^2 \right]_0^2= 
{4 \over 3}$$

Das Verhältnis der beiden Teilstücke ist: 
$$A_2= A_{ges}-A_1={9 \over 2}-{4\over 3}= {19 \over 6} ; \quad A_2 : A_1= 8 : 19$$

\noindent {\bf Aufgabe 6. d)}

Die zweite Ableitung der Integralfunktion $I(x)$ entspricht der ersten Ableitung von $f(x)$;
welche bereits oben diskutiert wurde. Es gilt: 
$$
\begin{array}{rll}
I''(x)>0 & \textnormal{für}~ x \in \left]-\infty ; 1,5\right] & I\ \textnormal{linksgekrümmt}  \\
I''(x)<0 & \textnormal{für}~ x \in \left[ 1,5;\infty \right[ & I\ \textnormal{rechtsgekrümmt}  \\
I''(x)=0 & \textnormal{für}~ x=1,5 & \Rightarrow \textnormal{Wendestelle}
\end{array}
$$

Die integralfreie Darstellung von $I(x)$ ist:

$$I(x)=\int \limits_0^x (-t^2+3t) {\rm d} t = \left[ -{1 \over 3} t^3 + {3\over 2} t^2 \right]_0^x= 
 -{1 \over 3} x^3 + {3\over 2} x^2$$

Die y-Koordinate des Wendepunktes von $I(x)$ ist also:

$$I(1,5)=-{1 \over 3} 1,5^3 + {3\over 2} 1,5^2=-1,125+3,375 = 2,25 \Rightarrow 
W(1,5|2,25)$$

\noindent{\bf Aufgabe 7.}\\
$$ 3 \ln (x^2+4) \qquad e^{x^2} \qquad {1\over 2} \sin (6x) $$

\noindent{\bf Aufgabe 8. a)}\\
Wir bilden die erste Ableitung und setzen sie gleich null:
$$
\renewcommand{\arraystretch}{2}
\begin{array}{rcll}
f'(x)&=&{3\over 4}x^2-6x+9=0 \\
    0&=& x^2 -8x +12 \\
  x_1&=& 6 \\
  x_2&=& 2
\end{array}
$$
Die zweite Ableitung ist:
$$
\renewcommand{\arraystretch}{2}
\begin{array}{rccll}
f''(x) &=& {3\over 2}x-6 \\
f''(x_1) &=& 9 -6 = 3&> 0 & \Rightarrow T (6|f(6))\\
f''(x_2) &=& 3 -6 = -3&< 0 & \Rightarrow H (2|f(2))\\
\end{array}
$$
Die Extrempunkte sind: $T(6|0)$ und $H(2|8)$.

\noindent{\bf Aufgabe 8 . b)}\\
Für die Wendestelle setzen wir die zweite Ableitung gleich null und erhalten:
$$f''(x) = {3\over 2}x-6 = 0 \quad \Rightarrow x_w =4 $$
Wir testen die dritte Ableitung an dieser Stelle: 
$$f'''(4)={3 \over 2} \not = 0$$
Also haben wir einen Wendepunkt bei $W(4|f(4))=(4|4)$ 

\noindent
\begin{minipage}{0.5\textwidth}
	\noindent{\bf Aufgabe 8. c)}\\
	Geradensteigung $m$:
	$$m = {\Delta y \over \Delta x} = {8-0 \over 2-6} = -2$$
	y-Achsenabschnitt $t$:
	$$ 0 = -2 \cdot 6 + t \quad \Rightarrow t = 12 $$
	Funktionsterm Gerade:
	$$ g(x) = -2 x + 12 $$
	Punktprobe Wendepunkt:
	$$ 4 = -2\cdot 4 + 12 \quad (w)$$
	
	
	Der Flächeninhalt $A$ zwischen $f(x)$ und der Geraden  setzt sich zusammen
	aus zwei Teilstücken $A_1, A_2$:
\end{minipage}
\hfill
\begin{minipage}{0.5\textwidth}
	\includegraphics[width=7 cm]{grint063}
\end{minipage}

$$A = |A_1|+|A_2|= |\int\limits_2^4 (f(x)-g(x)){\rm d} x | + |\int\limits_4^6 (f(x)-g(x)){\rm d} x | $$
Die Stammfunktion von $h(x)=f(x)-g(x)$ ist:
$$H(x)={1\over 16}x^4-x^3 +5,5x^2-12x +C $$
$$A =  |H(4)-H(2)|+|H(6)-H(4)|=|(-8)-(-9)|+|(-9)-(-8)|=2~{\rm FE} $$

\begin{center}
	\includegraphics[width=7 cm]{../../viecher/endcomic.pdf}
	
	Hier geht es zurück zum \href{https://www.okuyakl.de/math/m12ikdL063/aa063.pdf}{Aufgabenblatt}
\end{center}

\end{document}

